\section{Lie 群}\label{sec:lie-groups}

\textbf{Lie 群} $G$ 是同时也是群的光滑流形, 并且群运算
\[
  \begin{aligned}
    G\times G&\to G\\
    (g,h)&\mapsto gh
  \end{aligned}
  \quad\text{和}\quad
  \begin{aligned}
    G&\to G\\
    g&\mapsto g^{-1}
  \end{aligned}
\]
都是可微映射 (在 $G\times G$ 上取乘积的标准可微结构, 参见 Exercise~\ref{ex:exercise-differentiable-structures}(\ref{ex:product-differentiable-manifold})).

\begin{RiemannExample}\label{ex:classical-matrix-lie-groups}
\begin{enumerate}[label=(\arabic*),ref=\arabic*,leftmargin=*]
\item 显然, $(\mathbb{R}^n,+)$ 是 Abel Lie 群.
\item \textbf{一般线性群}
\[
  GL(n)=\{n\times n\text{ 可逆实矩阵}\}
\]
是最基本的非平凡 Lie 群的例子. 我们在 Example~\ref{ex:examples-of-differentiable-structures}(\ref{ex:matrix-space-and-general-linear-group}) 中已看到, 它是一个 $n^2$ 维光滑流形. 此外, 群乘法就是 $n\times n$ 矩阵的通常乘法到
\[
  GL(n)\times GL(n)
\]
的限制, 其坐标函数是二次多项式; 求逆就是非奇异矩阵的通常求逆到 $GL(n)$ 的限制. 由 Cramer 法则可知, 求逆的坐标函数是分母非零的有理函数 (分母中只出现行列式).
\item\label{ex:orthogonal-group-regular-level-set} $\mathbb{R}^n$ 的正交变换所成的\textbf{正交群}
\[
  O(n)=\{A\in\mathcal M_{n\times n}\mid A^tA=I\}
\]
也是 Lie 群. 为说明这一点, 考虑从 $\mathcal M_{n\times n}\cong\mathbb{R}^{n^2}$ 到 $n\times n$ 对称矩阵所成的空间 $\mathcal S_{n\times n}\cong\mathbb{R}^{\frac12 n(n+1)}$ 的映射 $f:A\mapsto A^tA$. 它在点 $A\in O(n)$ 处的导数 $(\mathrm{d}f)_A$ 是从 $T_A\mathcal M_{n\times n}\cong\mathcal M_{n\times n}$ 到 $T_{f(A)}\mathcal S_{n\times n}\cong\mathcal S_{n\times n}$ 的满射. 实际上,
\begin{align*}
  (\mathrm{d}f)_A(B)
  &=\lim_{h\to0}\frac{f(A+hB)-f(A)}{h}\\
  &=\lim_{h\to0}\frac{(A+hB)^t(A+hB)-A^tA}{h}\\
  &=B^tA+A^tB,
\end{align*}
且令 $B=\frac12(A^{-1})^tS=\frac12 AS$, 就能把任意对称矩阵 $S$ 写成 $B^tA+A^tB$. 特别地, 恒等矩阵 $I$ 是 $f$ 的正则值, 所以由 Theorem~\ref{thm:submanifold-of-regular-level-set} 可知, $O(n)=f^{-1}(I)$ 是 $\mathcal M_{n\times n}$ 的一个维数为 $\frac12 n(n-1)$ 的子流形. 此外, 它也是 Lie 群, 因为群乘法和求逆是 $GL(n)$ 上相应运算到子流形 $O(n)$ 的限制, 且它们的值仍在 $O(n)$ 中 [参见 Exercise~\ref{ex:immersions-embeddings}(\ref{ex:restriction-to-submanifolds})].
\item\label{ex:special-linear-group-regular-level-set} 由 $f(A)=\det A$ 给出的映射 $f:GL(n)\to\mathbb{R}$ 可微, 且水平集 $f^{-1}(1)$ 是\textbf{特殊线性群}
\[
  SL(n)=\{A\in\mathcal M_{n\times n}\mid\det A=1\}.
\]
同样, $f$ 在点 $A\in GL(n)$ 处的导数是满射, 这使得 $SL(n)$ 成为 Lie 群. 实际上, 容易看出
\[
  (\mathrm{d}f)_I(B)=\lim_{h\to0}\frac{\det(I+hB)-\det I}{h}
  =\operatorname{tr}B
\]
从而
\begin{align*}
  (\mathrm{d}f)_A(B)
  &=\lim_{h\to0}\frac{\det(A+hB)-\det A}{h}\\
  &=\lim_{h\to0}\frac{(\det A)\det(I+hA^{-1}B)-\det A}{h}\\
  &=(\det A)\lim_{h\to0}\frac{\det(I+hA^{-1}B)-1}{h}\\
  &=(\det A)(\mathrm{d}f)_I(A^{-1}B)
   =(\det A)\operatorname{tr}(A^{-1}B).
\end{align*}
由于 $\det(A)=1$, 对任意 $k\in\mathbb{R}$, 取矩阵 $B=\frac{k}{n}A$, 就有 $(\mathrm{d}f)_A(B)=\operatorname{tr}(\frac{k}{n}I)=k$. 因此, 对每个 $A\in SL(n)$, $(\mathrm{d}f)_A$ 都是满射, 所以 $1$ 是 $f$ 的一个正则值. 于是 $SL(n)$ 是 $GL(n)$ 的一个子流形. 和前一个例子一样, 群乘法和求逆可微, 所以 $SL(n)$ 是 Lie 群.
\item 映射 $A\mapsto\det A$ 是从 $O(n)$ 到 $\{-1,1\}$ 的可微映射, 水平集 $f^{-1}(1)$ 为
\[
  SO(n)=\{A\in O(n)\mid\det A=1\},
\]
称为 $\mathbb{R}^n$ 中的\textbf{特殊正交群}或\textbf{旋转群}. 它是 $O(n)$ 的开子集, 因而是同维数的 Lie 群.
\item 还可以考虑复 $n\times n$ 矩阵的空间 $\mathcal M_{n\times n}(\mathbb C)$ 和复可逆 $n\times n$ 矩阵的空间 $GL(n,\mathbb C)$. 后者是实维数为 $2n^2$ 的 Lie 群. 此外, 类似于上面对 $O(n)$ 的处理, 可以取 $\mathbb{C}^n$ 上的酉变换群
\[
  U(n)=\{A\in\mathcal M_{n\times n}(\mathbb{C})\mid A^*A=I\},
\]
其中 $A^*$ 是 $A$ 的伴随. 这个群是 $\mathcal M_{n\times n}(\mathbb{C})\cong\mathbb{C}^{n^2}\cong\mathbb{R}^{2n^2}$ 的子流形, 也是一个 Lie 群, 称为\textbf{酉群}. 这一点可由下述事实看出: $I$ 是从 $\mathcal M_{n\times n}(\mathbb{C})$ 到自伴矩阵空间的映射 $f:A\mapsto A^*A$ 的一个正则值. 任意 $\mathcal M_{n\times n}(\mathbb{C})$ 的元素都能唯一地写成一个自伴矩阵和一个反自伴矩阵之和, 且 $A\mapsto iA$ 是从自伴矩阵空间到反自伴矩阵空间的同构. 因此这两个空间的实维数均为 $\frac12\dim_{\mathbb{R}}\mathcal M_{n\times n}(\mathbb{C})=n^2$. 所以 $\dim U(n)=n^2$.
\item \textbf{特殊酉群}
\[
  SU(n)=\{A\in U(n)\mid\det A=1\},
\]
也是 Lie 群, 此时维数为 $n^2-1$ (注意, $A\mapsto\det(A)$ 现在是从 $U(n)$ 到 $S^1$ 的可微映射).
\end{enumerate}
\end{RiemannExample}

按定义, Lie 群 $G$ 是流形, 因而可以考虑它在某一点的切空间. 特别地, 在单位元 $e$ 处的切空间通常记为
\[
  \mathfrak g:=T_eG.
\]
对 $g\in G$, 有映射
\[
  \begin{aligned}
    L_g:G&\to G\\ h&\mapsto g\cdot h
  \end{aligned}
  \quad\text{和}\quad
  \begin{aligned}
    R_g:G&\to G\\ h&\mapsto h\cdot g
  \end{aligned}
\]
分别对应于 $g$ 的\textbf{左乘}和\textbf{右乘}.

若对每个 $g\in G$ 都有 $(L_g)_*X=X$, 则称 $G$ 上的向量场 $X$ \textbf{左不变}, 即对每个 $g,h\in G$,
\[
  ((L_g)_*X)_{gh}=X_{gh}\quad\text{或}\quad(\mathrm{d}L_g)_hX_h=X_{gh}.
\]
当然, $\mathfrak g$ 与 $G$ 上左不变向量场的空间之间存在向量空间同构, 它将每个 $V\in\mathfrak g$ 对应到由
\[
  X_g^V:=(\mathrm{d}L_g)_eV
\]
定义的向量场 $X^V$, 其中 $g\in G$ 任意. 这一向量场是左不变的, 因为
\begin{align*}
  (\mathrm{d}L_g)_hX_h^V
  &=(\mathrm{d}L_g)_h(\mathrm{d}L_h)_eV
    =(\mathrm{d}(L_g\circ L_h))_eV\\
  &=(\mathrm{d}L_{gh})_eV=X_{gh}^V.
\end{align*}
注意, 给定左不变向量场 $X$, $\mathfrak g$ 中相应的元素是 $X_e$. 左不变向量场的空间 $\mathfrak{X}_L(G)$ 在向量场的 Lie 括号下封闭 (因为由 Exercise~\ref{ex:vector-fields}(\ref{ex:pushforward-preserves-lie-bracket}), 有 $(L_g)_*[X,Y]=[(L_g)_*X,(L_g)_*Y]$), 因而它是向量场的 Lie 代数的一个 Lie 子代数 (见 Definition~\ref{def:lie-algebra}). 于是, 同构 $\mathfrak{X}_L(G)\cong\mathfrak g$ 在 $\mathfrak g$ 上确定了一个 Lie 代数结构. 我们称 $\mathfrak g$ 为 Lie 群 $G$ 的 \textbf{Lie 代数}.

\begin{RiemannExample}\label{ex:lie-algebras-of-matrix-groups}
\begin{enumerate}[label=(\arabic*),ref=\arabic*,leftmargin=*]
\item 若 $G=GL(n)$, 则 $\mathfrak{gl}(n)=T_I GL(n)=\mathcal M_{n\times n}$ 是实系数 $n\times n$ 矩阵的空间, $\mathfrak{gl}(n)$ 上的 Lie 括号是矩阵的交换子
\[
  [A,B]=AB-BA.
\]
实际上, 若 $A,B\in\mathfrak{gl}(n)$ 是两个 $n\times n$ 矩阵, 相应的左不变向量场由下式给出:
\begin{align*}
  X_g^A&=(\mathrm{d}L_g)_I(A)
       =\sum_{i,k,j}x^{ik}a^{kj}\frac{\partial}{\partial x^{ij}}\\
  X_g^B&=(\mathrm{d}L_g)_I(B)
       =\sum_{i,k,j}x^{ik}b^{kj}\frac{\partial}{\partial x^{ij}},
\end{align*}
其中 $g\in GL(n)$ 是分量为 $x^{ij}$ 的矩阵. $[X^A,X^B]_g$ 的 $ij$ 分量由 $X_g^A\cdot(X^B)^{ij}-X_g^B\cdot(X^A)^{ij}$ 给出, 即
\begin{align*}
  [X^A,X^B]^{ij}(g)
   &=\left(\sum_{l,m,p}x^{lp}a^{pm}\frac{\partial}{\partial x^{lm}}\right)
       \left(\sum_k x^{ik}b^{kj}\right)\\
   &\quad-\left(\sum_{l,m,p}x^{lp}b^{pm}\frac{\partial}{\partial x^{lm}}\right)
       \left(\sum_k x^{ik}a^{kj}\right)\\
   &=\sum_{k,l,m,p}x^{lp}a^{pm}\delta_{il}\delta_{km}b^{kj}
       -\sum_{k,l,m,p}x^{lp}b^{pm}\delta_{il}\delta_{km}a^{kj}\\
   &=\sum_{m,p}x^{ip}(a^{pm}b^{mj}-b^{pm}a^{mj})\\
   &=\sum_p x^{ip}(AB-BA)^{pj}
\end{align*}
(其中 $\delta_{ij}=1$ 若 $i=j$, $\delta_{ij}=0$ 若 $i\ne j$, 这是 Kronecker 符号). 令 $g=I$, 得
\[
  [A,B]=[X^A,X^B]_I=AB-BA.
\]
由 Exercise~\ref{ex:vector-fields}(\ref{ex:lie-bracket-tangent-to-submanifold}) 可知, 当 $G$ 是矩阵群, 即对某个 $n$, $G$ 是 $GL(n)$ 的子群时, 这一点总是成立.
\item\label{ex:lie-algebra-orthogonal-group} 若 $G=O(n)$, 则它的 Lie 代数为
\[
  \mathfrak o(n)=\{A\in\mathcal M_{n\times n}\mid A^t+A=0\}.
\]
实际上, 我们在 Example~\ref{ex:classical-matrix-lie-groups}(\ref{ex:orthogonal-group-regular-level-set}) 中已看到 $O(n)=f^{-1}(I)$, 其中恒等矩阵 $I$ 是映射
\begin{align*}
  f:\mathcal M_{n\times n}&\to\mathcal S_{n\times n}\\ A&\mapsto A^tA
\end{align*}
的正则值. 因此, $\mathfrak o(n)=T_IG=\ker(\mathrm{d}f)_I=\{A\in\mathcal M_{n\times n}\mid A^t+A=0\}$ 是反对称矩阵的空间.
\item 若 $G=SL(n)$, 则它的 Lie 代数为
\[
  \mathfrak{sl}(n)=\{A\in\mathcal M_{n\times n}\mid\operatorname{tr}A=0\}.
\]
实际上, 我们在 Example~\ref{ex:classical-matrix-lie-groups}(\ref{ex:special-linear-group-regular-level-set}) 中已看到 $SL(n)=f^{-1}(1)$, 其中 $1$ 是映射
\begin{align*}
  f:\mathcal M_{n\times n}&\to\mathbb{R}\\ A&\mapsto\det A
\end{align*}
的正则值. 因此, $\mathfrak{sl}(n)=T_IG=\ker(\mathrm{d}f)_I=\{A\in\mathcal M_{n\times n}\mid\operatorname{tr}A=0\}$ 是迹为零的矩阵的空间.
\item 若 $G=SO(n)=\{A\in O(n)\mid\det A=1\}$, 则它的 Lie 代数为
\[
  \mathfrak{so}(n)=T_I SO(n)=T_I O(n)=\mathfrak o(n).
\]
\item 类似于 Example~\ref{ex:lie-algebras-of-matrix-groups}(\ref{ex:lie-algebra-orthogonal-group}), $U(n)$ 的 Lie 代数为
\[
  \mathfrak u(n)=\{A\in\mathcal M_{n\times n}(\mathbb{C})\mid A^*+A=0\},
\]
即反 Hermite 矩阵的空间.
\item 为确定 $SU(n)$ 的 Lie 代数, 注意 $SU(n)$ 是 $f(A)=\det A$ 的水平集 $f^{-1}(1)$, 因而
\[
  \mathfrak{su}(n)=\ker(\mathrm{d}f)_I
  =\{A\in\mathfrak u(n)\mid\operatorname{tr}(A)=0\}.
\]
\end{enumerate}
\end{RiemannExample}

现在研究左不变向量场的流.

\begin{RiemannProposition}\label{prop:global-flow-of-left-invariant-vector-field}
设 $F$ 是左不变向量场 $X$ 在点 $h\in G$ 处的局部流. 则 $F$ 所定义的映射 $\psi_t$ (即 $\psi_t(g)=F(g,t)$) 满足 $\psi_t=R_{\psi_t(e)}$. 此外, $X$ 的流对所有 $t\in\mathbb{R}$ 全局有定义.
\end{RiemannProposition}

\begin{RiemannProof}
对 $g\in G$, 有 $R_{\psi_t(e)}(g)=g\cdot\psi_t(e)=L_g(\psi_t(e))$. 因此,
\[
  R_{\psi_0(e)}(g)=g\cdot e=g
\]
且
\begin{align*}
  \frac{\mathrm{d}}{\mathrm{d}t}\bigl(R_{\psi_t(e)}(g)\bigr)
    &=\frac{\mathrm{d}}{\mathrm{d}t}\bigl(L_g(\psi_t(e))\bigr)
      =(\mathrm{d}L_g)_{\psi_t(e)}\left(\frac{\mathrm{d}}{\mathrm{d}t}(\psi_t(e))\right)\\
    &=(\mathrm{d}L_g)_{\psi_t(e)}\bigl(X_{\psi_t(e)}\bigr)=X_{g\cdot\psi_t(e)}\\
    &=X_{R_{\psi_t(e)}(g)},
\end{align*}
因此 $R_{\psi_t(e)}(g)=c_g(t)=\psi_t(g)$ 是 $X$ 在 $g$ 处的积分曲线. 于是, 若 $\psi_t(e)$ 对 $t\in(-\varepsilon,\varepsilon)$ 有定义, 则 $\psi_t(g)$ 对 $t\in(-\varepsilon,\varepsilon)$, $g\in G$ 有定义. 此外, 第~\ref{sec:vector-fields}~节的条件 \eqref{eq:local-flow-group-law} 对每一组 $-\varepsilon/2<s,t<\varepsilon/2$ 都成立, 因此可以像前面那样把映射 $F$ 延拓到 $G\times\mathbb{R}$: 对任意 $t\in\mathbb{R}$, 写成 $t=k\varepsilon/2+s$, 其中 $k\in\mathbb{Z}$ 且 $0\leq s<\varepsilon/2$, 并定义 $F(g,t):=F^k(F(g,s),\varepsilon/2)=gF(e,s)F^k(e,\varepsilon/2)$.
\end{RiemannProof}

\begin{RiemannRemark}\label{rem:one-parameter-lie-group-homomorphism}
若 Lie 群之间的同态 $F:G_1\to G_2$ 除了是群同态之外, 还是可微映射, 就称为 \textbf{Lie 群同态}. 由于
\[
  \psi_{t+s}(e)=\psi_s(\psi_t(e))=R_{\psi_s(e)}\psi_t(e)=\psi_t(e)\cdot\psi_s(e),
\]
积分曲线 $t\mapsto\psi_t(e)$ 定义了 $(\mathbb{R},+)$ 到 $(G,\cdot)$ 的群同态.
\end{RiemannRemark}

\begin{RiemannDefinition}\label{def:lie-group-exponential-map}
\textbf{指数映射} $\exp:\mathfrak g\to G$ 将每个 $V\in\mathfrak g$ 映到 $\psi_1(e)$, 其中 $\psi_t$ 是左不变向量场 $X^V$ 的流.
\end{RiemannDefinition}

\begin{RiemannRemark}\label{rem:one-parameter-subgroup-via-exponential}
若 $c_g(t)$ 是 $X$ 在 $g$ 处的积分曲线且 $s\in\mathbb{R}$, 容易验证 $c_g(st)$ 是 $sX$ 在 $g$ 处的积分曲线. 另一方面, 对 $V\in\mathfrak g$ 有 $X^{sV}=sX^V$. 因此,
\[
  \psi_t(e)=c_e(t)=c_e(t\cdot 1)=F(e,1)=\exp(tV),
\]
其中 $F$ 是 $tX^V=X^{tV}$ 的流.
\end{RiemannRemark}

\begin{RiemannExample}\label{ex:matrix-lie-group-exponential}
若 $G$ 是矩阵群, 则对 $A\in\mathfrak g$,
\[
  \exp A=e^A=\sum_{k=0}^{\infty}\frac{A^k}{k!}.
\]
实际上, 对任意矩阵 $A$, 这一系列都收敛, 而且映射 $h(t)=e^{At}$ 满足
\begin{align*}
  h(0)&=e^0=I\\
  \frac{\mathrm{d}h}{\mathrm{d}t}(t)&=e^{At}A=h(t)A.
\end{align*}
因此, $h$ 是 $X^A$ 在单位元处的流 (即 $h(t)=\psi_t(e)$), 故 $\exp A=\psi_1(e)=e^A$.
\end{RiemannExample}

现在设 $G$ 是任意群, $M$ 是任意集合. 若存在从 $G$ 到 $M$ 到自身的双射映射所成的群的同态 $\phi$, 就称 \phantomsection\label{idx:group-action-definition}\textbf{$G$ 作用于 $M$}. 等价地, 写作
\[
  \phi(g)(p)=A(g,p),
\]
则要求存在满足以下条件的映射 $A:G\times M\to M$:
\begin{enumerate}[label=(\roman*),leftmargin=*]
\item 若 $e$ 是 $G$ 的单位元, 则 $A(e,p)=p$, $\forall p\in M$;
\item 若 $g,h\in G$, 则 $A(g,A(h,p))=A(gh,p)$, $\forall p\in M$.
\end{enumerate}
通常用 $g\cdot p$ 表示 $A(g,p)$.

\begin{RiemannExample}\label{ex:group-actions-by-left-multiplication}
\begin{enumerate}[label=(\arabic*),leftmargin=*]
\item 设 $G$ 是群, $H\subset G$ 是子群. 则 $H$ 通过左乘作用于 $G$: 对 $h\in H$, $g\in G$, 有 $A(h,g)=h\cdot g$.
\item $GL(n)$ 通过 $A\cdot x=Ax$ 作用于 $\mathbb{R}^n$, 其中 $A\in GL(n)$ 且 $x\in\mathbb{R}^n$. 对任何子群 $G\subset GL(n)$ 也是如此.
\end{enumerate}
\end{RiemannExample}

对每个 $p\in M$, 可以把 $p$ 的\textbf{轨道}定义为集合 $G\cdot p:=\{g\cdot p\mid g\in G\}$. 若 $G\cdot p=\{p\}$, 则称 $p$ 为 $G$ 的\textbf{不动点}. 若存在点 $p\in M$ 的轨道是整个 $M$ (即 $G\cdot p=M$), 则称作用是\textbf{传递的}. 注意, 此时只有一条轨道, 且对任意 $p,q\in M$ 且 $p\ne q$, 总存在群的元素 $g\in G$ 使得 $q=g\cdot p$. 这时称流形 $M$ 为 $G$ 的\textbf{齐性空间}. 点 $p\in M$ 的\textbf{稳定子} (或\textbf{各向同性子群}) 是群
\[
  G_p=\{g\in G\mid g\cdot p=p\}.
\]
若所有稳定子都是平凡的, 则称作用是\textbf{自由的}.

若 $G$ 是 Lie 群, $M$ 是光滑流形, 当映射 $A:G\times M\to M$ 可微时, 称作用是\textbf{光滑的}. 此时, 映射 $p\mapsto g\cdot p$ 是微分同胚. 我们始终假设 Lie 群在可微流形上的作用是光滑的. 若映射
\begin{align*}
  G\times M&\to M\times M\\
  (g,p)&\mapsto(g\cdot p,p)
\end{align*}
是固有映射, 则称这个光滑作用是\textbf{固有的} (回顾一下, 若每个紧集的原像都紧, 则称映射是\textbf{固有映射}, 参见第~\RGref{subsec:notes-proper-maps-and-fundamental-group}~节).

\begin{RiemannRemark}\label{rem:sequential-characterization-of-proper-actions}
注意, 一个光滑作用是固有的, 当且仅当给定 $M$ 中两个收敛序列 $\{p_n\}$ 和 $\{g_n\cdot p_n\}$, $G$ 中存在收敛子序列 $\{g_{n_k}\}$. 若 $G$ 紧致, 这个条件总是成立.
\end{RiemannRemark}

$G$ 在 $M$ 上的作用的轨道是由 $p\sim q\Leftrightarrow q\in G\cdot p$ 给出的等价关系 $\sim$ 的等价类 (参见第~\RGref{subsec:notes-topological-manifolds}~节). 因此, 商 (拓扑) 空间 $M/{\sim}$ 通常称为该作用的\textbf{轨道空间}, 并记作 $M/G$.

\begin{RiemannProposition}\label{prop:hausdorff-orbit-space}
若 Lie 群 $G$ 在可微流形 $M$ 上的作用是固有的, 则轨道空间 $M/G$ 是 Hausdorff 空间.
\end{RiemannProposition}

\begin{RiemannProof}
关系 $p\sim q\Leftrightarrow q\in G\cdot p$ 是开等价关系 (参见第~\RGref{subsec:notes-topological-manifolds}~节). 实际上, 由于 $p\mapsto g\cdot p$ 是同胚, 对 $M$ 中的任何开集 $U$, 集合 $\pi^{-1}(\pi(U))=\{g\cdot p\mid p\in U\text{ 且 }g\in G\}=\bigcup_{g\in G}g\cdot U$ 是 $M$ 的开子集, 即 $\pi(U)$ 是开集 (这里 $\pi:M\to M/G$ 是商映射). 因此只需证明集合
\[
  R=\{(p,q)\in M\times M\mid p\sim q\}
\]
是闭的 (参见 Proposition~\RGref{prop:topology-of-quotients}). 这一点由 $R$ 是映射
\begin{align*}
  G\times M&\to M\times M\\
  (g,p)&\mapsto(g\cdot p,p)
\end{align*}
的像推出, 因为该映射连续而且固有, 因而是闭映射 (参见第~\RGref{subsec:notes-proper-maps-and-fundamental-group}~节).
\end{RiemannProof}

在一定条件下, 轨道空间 $M/G$ 自然成为可微流形.

\begin{RiemannTheorem}\label{thm:quotient-manifold-of-free-proper-action}
设 $M$ 是可微流形, 赋有 Lie 群 $G$ 的一个自由且固有的作用. 则轨道空间 $M/G$ 自然成为维数为 $\dim M-\dim G$ 的可微流形, 而且商映射 $\pi:M\to M/G$ 是淹没.
\end{RiemannTheorem}

\begin{RiemannProof}
由前一命题, 商空间 $M/G$ 是 Hausdorff 的. 此外, 由于 $M$ 满足第二可数公理, 且 $G$ 所定义的等价关系是开等价关系, 这个商空间也满足第二可数公理. 还需证明 $M/G$ 具有一种自然的可微结构, 使商映射是淹没. 我们只证明 $G$ 是离散 (即零维) Lie 群的情形 (参见 Remark~\ref{rem:topological-manifolds-low-dimension}); 一般 Lie 群的证明可见 \RGcite{DK99}.

在这里, 只需证明: 对每个 $p\in M$, 存在邻域 $U\ni p$, 使得当 $g\ne h$ 时, $g\cdot U\cap h\cdot U=\varnothing$. 这保证每个点 $[p]\in M/G$ 都有一个与 $U$ 同胚的邻域 $[U]$, 并可假定 $U$ 是坐标邻域. 由于 $G$ 通过微分同胚作用, 如此定义的可微结构与 $p\in[p]$ 的选取无关. $M/G$ 的坐标图来自 $M$ 的坐标图, 所以重叠映射是光滑的. 因此, $M/G$ 具有一种自然的可微结构, 使 $\pi:M\to M/G$ 是局部微分同胚 (因为 $\pi|_U:U\to[U]$ 的坐标表达式是恒等映射).

证明当 $g\ne h$ 时 $g\cdot U\cap h\cdot U=\varnothing$, 等价于证明当 $g\ne e$ 时 $g\cdot U\cap U=\varnothing$. 假设对于任何邻域 $U\ni p$, 这一点都不成立. 则存在一列开集 $U_n\ni p$, 满足 $U_{n+1}\subset U_n$, $\bigcap_{n=1}^{+\infty}U_n=\{p\}$, 以及一列 $g_n\in G\setminus\{e\}$, 使得 $g_n\cdot U_n\cap U_n\ne\varnothing$. 取 $p_n\in g_n\cdot U_n\cap U_n$. 则对某个 $q_n\in U_n$, 有 $p_n=g_n\cdot q_n$. 我们有 $p_n\to p$ 及 $q_n\to p$. 由于作用固有, $g_n$ 存在收敛子序列 $g_{n_k}$. 设 $g$ 为其极限. 在 $q_{n_k}=g_{n_k}\cdot p_{n_k}$ 中令 $k\to+\infty$, 得 $g\cdot p=p$, 从而 $g=e$ (因为作用是自由的). 由于 $G$ 离散, 对充分大的 $k$ 应有 $g_{n_k}=e$, 矛盾.
\end{RiemannProof}

\begin{RiemannExample}\label{ex:projective-spaces-and-tori-as-quotients}
\begin{enumerate}[label=(\arabic*),ref=\arabic*,leftmargin=*]
\item\label{ex:projective-space-antipodal-quotient} 设 $S^n=\{x\in\mathbb{R}^{n+1}\mid\sum_{i=1}^{n}(x^i)^2=1\}$ 赋有 $G=\mathbb{Z}_2=\{-I,I\}$ 由 $-I\cdot x=-x$ (对径映射) 给出的作用. 这个作用既固有又自由, 所以轨道空间 $S^n/G$ 是 $n$ 维流形. 这个空间就是实射影空间 $\mathbb{R}P^n$ [参见 Exercise~\ref{ex:exercise-differentiable-structures}(\ref{ex:projective-space-atlas-from-sphere})].
\item 群 $G=\mathbb{R}\setminus\{0\}$ 通过乘法 $t\cdot x=tx$ 作用于 $M=\mathbb{R}^{n+1}\setminus\{0\}$. 这个作用既固有又自由, 因而 $M/G$ 是 $n$ 维可微流形 (它又是 $\mathbb{R}P^n$).
\item 考虑 $M=\mathbb{R}^n$, 赋有 $G=\mathbb{Z}^n$ 由下式定义的作用:
\[
  (k^1,\ldots,k^n)\cdot(x^1,\ldots,x^n)
  =(x^1+k^1,\ldots,x^n+k^n).
\]
这个作用既固有又自由, 所以商空间 $M/G$ 是 $n$ 维流形. 赋予 Theorem~\ref{thm:quotient-manifold-of-free-proper-action} 所定义的商可微结构后, 这个空间称为 \textbf{$n$ 维环面}, 记为 $T^n$. 它与乘积流形 $S^1\times\cdots\times S^1$ 微分同胚, 且当 $n=2$ 时, 它与 $\mathbb{R}^3$ 中的旋转环面微分同胚.
\end{enumerate}
\end{RiemannExample}

离散群作用的商给出了流形的\textbf{覆盖}.

\begin{RiemannDefinition}\label{def:smooth-covering-of-manifold}
可微流形 $B$ 的一个\textbf{光滑覆盖}是二元组 $(M,\pi)$, 其中 $M$ 是连通可微流形, $\pi:M\to B$ 是满射的局部微分同胚, 且对每个 $p\in B$, 在 $B$ 中都存在 $p$ 的连通邻域 $U$, 使得 $\pi^{-1}(U)$ 是互不相交的开集 $U_\alpha\subset M$ (称为\textbf{片}) 的并, 且 $\pi$ 到 $U_\alpha$ 的限制 $\pi_\alpha$ 都是到 $U$ 的微分同胚. 映射 $\pi$ 称为\textbf{覆盖映射}, $M$ 称为\textbf{覆盖流形}.
\end{RiemannDefinition}

\begin{RiemannRemark}\label{rem:dimensions-and-fibers-of-coverings}
\begin{enumerate}[label=(\arabic*),leftmargin=*]
\item 显然, 必有 $\dim M=\dim B$.
\item 注意, 互不相交的开集族 $\{U_\alpha\}$ 必为可数 ($M$ 有可数基).
\item \textbf{纤维} $\pi^{-1}(p)\subset M$ 具有离散拓扑. 实际上, 每一片 $U_\alpha$ 都是开集且与 $\pi^{-1}(p)$ 恰好相交于一点, 因此该点在子空间拓扑中是开集.
\end{enumerate}
\end{RiemannRemark}

\begin{RiemannExample}\label{ex:coverings-of-circle-torus-projective-space}
\begin{enumerate}[label=(\arabic*),ref=\arabic*,leftmargin=*]
\item\label{ex:real-line-covers-circle} 由
\[
  \pi(t)=(\cos(2\pi t),\sin(2\pi t))
\]
给出的映射 $\pi:\mathbb{R}\to S^1$ 是 $S^1$ 的一个光滑覆盖. 然而, 这个映射限制到 $(0,+\infty)$ 后是满射的局部微分同胚, 却不是覆盖映射.
\item\label{ex:euclidean-space-covers-torus} 覆盖映射的乘积显然是覆盖映射. 因此, 可以推广上例, 得到以 $\mathbb{R}^n$ 覆盖 $T^n\cong S^1\times\cdots\times S^1$ 的覆盖.
\item\label{ex:sphere-covers-projective-space} 在 Example~\ref{ex:projective-spaces-and-tori-as-quotients}(\ref{ex:projective-space-antipodal-quotient}) 中, 得到了以 $S^n$ 覆盖 $\mathbb{R}P^n$ 的覆盖.
\end{enumerate}
\end{RiemannExample}

设 $M$ 是覆盖流形. 若微分同胚 $h:M\to M$ 满足 $\pi\circ h=\pi$, 或等价地, $h$ 将每个集合 $\pi^{-1}(p)$ 映到自身, 则称 $h$ 为\textbf{甲板变换} (或\textbf{覆盖变换}). 可以证明, 全体覆盖变换所成的群 $G$ 是离散 Lie 群, 它在 $M$ 上的作用既自由又固有.

若覆盖流形 $M$ 是\textbf{单连通的} (参见第~\RGref{subsec:notes-proper-maps-and-fundamental-group}~节), 则称这个覆盖为\textbf{万有覆盖}. 此时, $B$ 与 $M/G$ 微分同胚. 此外, $G$ 同构于 $B$ 的\textbf{基本群} $\pi_1(B)$ (参见第~\RGref{subsec:notes-proper-maps-and-fundamental-group}~节).

\textbf{Lie 定理} (例如见 \RGcite{DK99}) 断言, 对给定的 Lie 代数 $\mathfrak g$, 存在唯一的单连通 Lie 群 $\widetilde G$, 其 Lie 代数为 $\mathfrak g$. 若 Lie 群 $G$ 的 Lie 代数也是 $\mathfrak g$, 则存在唯一的 Lie 群同态 $\pi:\widetilde G\to G$, 它是覆盖映射. 此时甲板变换群就是 $\ker(\pi)$, 因而 $G$ 与 $\widetilde G/\ker(\pi)$ 微分同胚. 实际上, $G$ 也同构于 $\widetilde G/\ker(\pi)$; 后者具有自然的群结构 (因为 $\ker(\pi)$ 是正规子群).

\begin{RiemannExample}\label{ex:fundamental-groups-from-universal-coverings}
\begin{enumerate}[label=(\arabic*),leftmargin=*]
\item 在 Example~\ref{ex:coverings-of-circle-torus-projective-space}(\ref{ex:real-line-covers-circle}) 的 $S^1$ 的万有覆盖中, 甲板变换是整数 $k$ 所给的平移 $h_k:t\mapsto t+k$, 因而 $S^1$ 的基本群是 $\mathbb{Z}$.
\item 类似地, $T^n$ 的万有覆盖的甲板变换是按整数向量的平移 [参见 Example~\ref{ex:coverings-of-circle-torus-projective-space}(\ref{ex:euclidean-space-covers-torus})], 因而 $T^n$ 的基本群是 $\mathbb{Z}^n$.
\item 在 Example~\ref{ex:coverings-of-circle-torus-projective-space}(\ref{ex:sphere-covers-projective-space}) 的 $\mathbb{R}P^n$ 的万有覆盖中, 仅有的甲板变换是恒等映射和对径映射, 因而 $\mathbb{R}P^n$ 的基本群是 $\mathbb{Z}_2$.
\end{enumerate}
\end{RiemannExample}

\begin{RiemannExercise}\label{ex:lie-groups}
\begin{enumerate}[label=(\arabic*),ref=\arabic*,leftmargin=*]
\item\label{ex:products-of-lie-groups-and-tori}
\begin{enumerate}[label=(\alph*),leftmargin=*]
\item 给定两个 Lie 群 $G_1,G_2$, 证明 $G_1\times G_2$ (两个群的直积) 赋有乘积的标准可微结构后是 Lie 群.
\item 圆 $S^1$ 可以与绝对值为 $1$ 的复数集合等同. 证明 $S^1$ 是 Lie 群, 并由此得出 $n$ 维环面 $T^n\cong S^1\times\cdots\times S^1$ 也是 Lie 群.
\end{enumerate}
\item\label{ex:abelian-lie-groups-and-exponential}
\begin{enumerate}[label=(\alph*),leftmargin=*]
\item 证明 $(\mathbb{R}^n,+)$ 是 Lie 群, 确定其 Lie 代数并写出指数映射的表达式.
\item 证明若 $G$ 是 Abel Lie 群, 则对所有 $V,W\in\mathfrak g$, 有 $[V,W]=0$.
\end{enumerate}
\item\label{ex:affine-maps-lie-group} 可以将
\[
  H=\{(x,y)\in\mathbb{R}^2\mid y>0\}
\]
中的每个点与由 $h(t)=yt+x$ 给出的可逆仿射映射 $h:\mathbb{R}\to\mathbb{R}$ 等同. 这样的映射全体关于复合构成一个群, 因而这一等同在 $H$ 上诱导出一个群结构.
\begin{enumerate}[label=(\alph*),leftmargin=*]
\item 证明诱导出的群运算由
\[
  (x,y)\cdot(z,w)=(yz+x,yw),
\]
给出, 且 $H$ 赋予这一群运算后是 Lie 群.
\item 证明左平移映射 $L_{(x,y)}:H\to H$ 在点 $(z,w)\in H$ 处的导数, 在上述坐标下由矩阵
\[
  (\mathrm{d}L_{(x,y)})_{(z,w)}=\begin{pmatrix}y&0\\0&y\end{pmatrix}.
\]
表示. 由此得出, 由向量
\[
  V=\xi\frac{\partial}{\partial x}+\eta\frac{\partial}{\partial y}
   \in\mathfrak h\equiv T_{(0,1)}H\qquad(\xi,\eta\in\mathbb{R})
\]
确定的左不变向量场 $X^V\in\mathfrak{X}(H)$ 由下式给出:
\[
  X_{(x,y)}^V=\xi y\frac{\partial}{\partial x}+\eta y\frac{\partial}{\partial y}.
\]
\item 给定 $V,W\in\mathfrak h$, 计算 $[V,W]$.
\item 确定向量场 $X^V$ 的流, 并给出指数映射 $\exp:\mathfrak h\to H$ 的表达式.
\item 先证明 $H$ 是 $GL(2)$ 中由满足 $y>0$ 的矩阵
\[
  \begin{pmatrix}y&x\\0&1\end{pmatrix}
\]
组成的子群, 再验证你的结果.
\end{enumerate}
\item\label{ex:special-linear-two-diffeomorphism} 考虑群
\[
  SL(2)=\left\{\begin{pmatrix}a&b\\c&d\end{pmatrix}\mid ad-bc=1\right\}.
\]
我们已经知道它是 3 维流形. 令
\[
  a=p+q,\quad d=p-q,\quad b=r+s,\quad c=r-s,
\]
证明 $SL(2)$ 与 $S^1\times\mathbb{R}^2$ 微分同胚.
\item\label{ex:exponentials-of-noncommuting-matrices} 举出矩阵 $A,B\in\mathfrak{gl}(2)$ 的例子, 使得 $e^{A+B}\ne e^Ae^B$.
\item\label{ex:determinant-of-matrix-exponential} 对 $A\in\mathfrak{gl}(n)$, 考虑可微映射
\begin{align*}
  h:\mathbb{R}&\to\mathbb{R}\setminus\{0\}\\
  t&\mapsto\det e^{At}
\end{align*}
并证明:
\begin{enumerate}[label=(\alph*),leftmargin=*]
\item 该映射是 $(\mathbb{R},+)$ 到 $(\mathbb{R}\setminus\{0\},\cdot)$ 的群同态;
\item $h'(0)=\operatorname{tr}A$;
\item $\det(e^A)=e^{\operatorname{tr}A}$.
\end{enumerate}
\item\label{ex:non-surjective-special-linear-exponential}
\begin{enumerate}[label=(\alph*),leftmargin=*]
\item 若 $A\in\mathfrak{sl}(2)$, 证明存在 $\lambda\in\mathbb{R}\cup i\mathbb{R}$, 使得
\[
  e^A=\cosh\lambda\,I+\frac{\sinh\lambda}{\lambda}A.
\]
\item 证明 $\exp:\mathfrak{sl}(2)\to SL(2)$ 不是满射.
\end{enumerate}
\item\label{ex:free-action-on-punctured-plane} 考虑由下式定义的向量场 $X\in\mathfrak{X}(\mathbb{R}^2)$:
\[
  X=\sqrt{x^2+y^2}\,\frac{\partial}{\partial x}.
\]
\begin{enumerate}[label=(\alph*),leftmargin=*]
\item 证明 $X$ 的流定义了 $\mathbb{R}$ 在 $M=\mathbb{R}^2\setminus\{0\}$ 上的自由作用.
\item 描述商拓扑空间 $M/\mathbb{R}$. 上述作用是固有的吗?
\end{enumerate}
\item\label{ex:diagonal-circle-action-on-product-of-spheres} 设 $M=S^2\times S^2$, 考虑 $S^1$ 在 $M$ 上由下式给出的对角作用:
\[
  e^{i\theta}\cdot(u,v)=\bigl(e^{i\theta}\cdot u,e^{2i\theta}\cdot v\bigr),
\]
其中, 对 $u\in S^2\subset\mathbb{R}^3$ 及 $e^{i\beta}\in S^1$, $e^{i\beta}\cdot u$ 表示将 $u$ 绕 $z$ 轴旋转角度 $\beta$.
\begin{enumerate}[label=(\alph*),leftmargin=*]
\item 确定此作用的不动点.
\item 可能的非平凡稳定子有哪些?
\end{enumerate}
\item\label{ex:closed-lie-subgroup-action-is-proper} 设 $G$ 是 Lie 群, $H$ 是闭 Lie 子群, 即 $H$ 是 $G$ 的子群且也是 $G$ 的闭子流形. 证明 $H$ 在 $G$ 上由 $A(h,g)=h\cdot g$ 定义的作用既自由又固有.
\item\label{ex:grassmannian-as-a-quotient} (\emph{Grassmann 流形}) 考虑形如
\[
  \begin{pmatrix}A&0\\C&B\end{pmatrix}
\]
的可逆矩阵所成的集合 $H\subset GL(n)$, 其中 $A\in GL(k)$, $B\in GL(n-k)$, $C\in\mathcal M_{(n-k)\times k}$.
\begin{enumerate}[label=(\alph*),leftmargin=*]
\item 证明 $H$ 是 $GL(n)$ 的闭 Lie 子群. 因而 $H$ 在 $GL(n)$ 上的作用既自由又固有 [参见 Exercise~\ref{ex:lie-groups}(\ref{ex:closed-lie-subgroup-action-is-proper})].
\item 证明商流形
\[
  Gr(n,k):=GL(n)/H
\]
可以与 $\mathbb{R}^n$ 的 $k$ 维子空间所成的集合等同 (特别地, $Gr(n,1)$ 就是射影空间 $\mathbb{R}P^{n-1}$).
\item 流形 $Gr(n,k)$ 称为 $\mathbb{R}^n$ 中 $k$ 维平面的 \textbf{Grassmann 流形}. 它的维数是多少?
\end{enumerate}
\item\label{ex:differential-of-lie-group-homomorphism} 设 $G$ 与 $H$ 是 Lie 群, $F:G\to H$ 是 Lie 群同态. 证明:
\begin{enumerate}[label=(\alph*),leftmargin=*]
\item $(\mathrm{d}F)_e:\mathfrak g\to\mathfrak h$ 是 Lie 代数同态;
\item 若 $(\mathrm{d}F)_e$ 是同构, 则 $F$ 是局部微分同胚;
\item 若 $F$ 是满射的局部微分同胚, 则 $F$ 是覆盖映射.
\end{enumerate}
\item\label{ex:quaternions-and-covering-of-rotation-group}
\begin{enumerate}[label=(\alph*),leftmargin=*]
\item 证明 $\mathbb{R}\cdot SU(2)$ 是 $\mathcal M_{2\times 2}(\mathbb{C})$ 的 4 维实线性子空间, 在矩阵乘法下封闭, 并有一组基
\begin{align*}
  1&=\begin{pmatrix}1&0\\0&1\end{pmatrix},&
  i&=\begin{pmatrix}i&0\\0&-i\end{pmatrix},\\
  j&=\begin{pmatrix}0&1\\-1&0\end{pmatrix},&
  k&=\begin{pmatrix}0&i\\i&0\end{pmatrix},
\end{align*}
满足 $i^2=j^2=k^2=ijk=-1$. 因此, 这一空间可以与\textbf{四元数}等同 (参见第~\RGref{subsec:notes-topological-manifolds}~节). 证明 $SU(2)$ 可以与 Euclidean 范数等于 $1$ 的四元数等同, 因而与 $S^3$ 微分同胚.
\item 若 $n\in\mathbb{R}^3$ 是一个单位向量, 将其与实部为零的一个四元数等同, 证明
\[
  \exp\left(\frac{n\theta}{2}\right)
    =1\cos\left(\frac{\theta}{2}\right)+n\sin\left(\frac{\theta}{2}\right)
\]
也是单位四元数.
\item 再将 $\mathbb{R}^3$ 与实部为零的四元数等同, 证明映射
\begin{align*}
  \mathbb{R}^3&\to\mathbb{R}^3\\
  v&\mapsto\exp\left(\frac{n\theta}{2}\right)\cdot v\cdot
       \exp\left(-\frac{n\theta}{2}\right)
\end{align*}
是绕 $n$ 确定的轴旋转角度 $\theta$ 的旋转.
\item 证明存在满射同态 $F:SU(2)\to SO(3)$, 并借此得出 $SU(2)$ 是 $SO(3)$ 的万有覆盖.
\item $SO(3)$ 的基本群是什么?
\end{enumerate}
\end{enumerate}
\end{RiemannExercise}

